\section{Bubonic plague}
Yersinia pestis (Y. pestis) was first identified in 1890 by Alexander Yersin. It was later attributed (and only recently confirmed with genomic data) to have caused a number of major epidemics which killed a huge proportion of the European population and went on to shape the course of western cultural development. 


In 751AD there was Justinian's plague, but the most dramatic epidemic started in 1345 and may have killed as much as 50\% of the European population \cite{belich2022world}, with this figure reaching above 80\% in some urban centres. Successive waves of infection followed this and the continent didn't fully emerge from the plight until as late as the 17th century, with the 1666 fire of London often attributed as the watershed moment.\par
There were likely earlier Y. pestis epidemics reaching far back into history beyond the scope of accurate written account, perhaps recorded only in legend and myth. Evidence for this can be seen in a recent discovery that Y. pestis emerged from its evolutionary antecedent (Yersinia pseudotuberculosis) around 6000 years ago during what would have been a period of human population expansion driven by the development of agriculture. We may hypothesise that it was this excess supply of replicative means which provided the raw material for evolution. Facilitating a shift in phenotypes over the saddle into a new local optima. This representing the niche pathogenesis observed today in Y. pestis.
